\documentclass[11pt,a4paper]{article}
\usepackage{gtu-el-yazisi}
\begin{document}
\notbilgi{Sayısal Haberleşme}{Örnekleme teoremi ve örnekleme türleri}{Ders-4.pdf}

\sayfa{1}
\elcizimi[0.9]{sekiller/Ders-4-s1-1.png}[1]

\section{Dirac katarı}
\[ x_\delta(t)=\sum_{n=-\infty}^{+\infty}\delta(t-nT_s) \;\overset{F}{\longleftrightarrow}\; X_\delta(f)=\overset{=f_s}{\kutula{\frac{1}{T_s}}}\sum_{n=-\infty}^{+\infty}\delta(f-nf_s) \]

\elcizimi[0.9]{sekiller/Ders-4-s1-2.png}[1]

\begin{align*}
x_s(t) &= x(t)\cdot x_\delta(t)=x(t)\sum_{n=-\infty}^{\infty}\delta(t-nT_s)\\
&=\sum_{n=-\infty}^{+\infty}x(nT_s)\,\delta(t-nT_s) \qquad \textcolor{turuncukalem}{\bigl(f(t)\,\delta(t-t_0)=f(t_0)\,\delta(t-t_0)\bigr)}
\end{align*}
\noindent $\uparrow$ Örneklenmiş işaret (Dirac katarı ile)

\elcizimi[0.9]{sekiller/Ders-4-s1-3.png}[1]

\noindent $f_s-f_m>f_m$ \\ $f_s>2f_m$

\noindent $f_s\ge(f_N=2f_m)$
\begin{align*}
x_s(t) &= x(t)\cdot x_\delta(t)=x(t)\sum_{n=-\infty}^{\infty}\delta(t-nT_s)\\
&= \kutula{\sum_{n=-\infty}^{\infty}x(nT_s)\,\delta(t-nT_s)}
\end{align*}
\formulkutu{f(t)\,\delta(t-t_0)=f(t_0)\,\delta(t-t_0)}
\[ \updownarrow F \]
\begin{align*}
X_s(f) &= X(f)\otimes X_\delta(f)\\
&= X(f)\otimes f_s\sum_{n=-\infty}^{\infty}\delta(f-nf_s)\\
&= f_s\sum_{n=-\infty}^{+\infty}X(f)\otimes\delta(f-nf_s)
\end{align*}
\[ \kutula{X_s(f)=f_s\sum_{n=-\infty}^{\infty}X(f-nf_s)} \]
\formulkutu{X(f)\otimes\delta(f-f_0)=X(f-f_0)}
\noindent $\uparrow$ sbt

\noindent $f_s<2f_m$ olursa \qquad Örtüşme (Aliasing)

\elcizimi[0.5]{sekiller/Ders-4-s1-4.png}[1]

\sayfa{2}
\noindent Pratikte $x(t)$ zamanda sınırlı, frekansta sınırlı değil (not band limited)

\elcizimi[0.9]{sekiller/Ders-4-s2-1.png}[2]

\noindent Çözüm: Örtüşme engelleyici filtre \\ Anti-aliasing filter

\elcizimi[0.6]{sekiller/Ders-4-s2-2.png}[2]

\section{İşareti yeniden oluşturma}
\noindent (Signal reconstruction)

\elcizimi[0.45]{sekiller/Ders-4-s2-3.png}[2]

\[ H(f)=A\,\Pi\!\left(\frac{f}{T}\right)=T_s\,\Pi\!\left(\frac{f}{2f_m}\right) \]
\[ \downarrow F^{-1} \]
\formulkutu[önemli]{h(t)=2f_mT_s\,\mathrm{sinc}(2f_mt)}
\noindent Eğer $f_s=2f_m \;\rightarrow\;$ 1 olur

\formulkutu{\begin{array}{c} A\,\Pi\!\left(\frac{t}{T}\right)\overset{F}{\longleftrightarrow}AT\,\mathrm{sinc}(fT)\\[4pt] AT\,\mathrm{sinc}(tT)\overset{F}{\longleftrightarrow}A\,\Pi\!\left(\frac{f}{T}\right)\end{array}}

\begin{align*}
\hat{x}(t) &= x_s(t)\otimes h(t)\\
&= h(t)\otimes x_s(t)=h(t)\otimes\sum_{n=-\infty}^{\infty}x(nT_s)\,\delta(t-nT_s)\\
&= \sum_{n=-\infty}^{\infty}x(nT_s)\,\underbrace{\bigl[h(t)\otimes\delta(t-nT_s)\bigr]}_{\textcolor{kirmizikalem}{h(t-nT_s)}}
\end{align*}
\begin{align*}
\hat{x}(t) &= \sum_{n=-\infty}^{\infty}x(nT_s)\,h(t-nT_s)\\
&= 2f_mT_s\sum_{n=-\infty}^{\infty}x(nT_s)\,\mathrm{sinc}\bigl[2f_m(t-nT_s)\bigr]
\end{align*}
\noindent Üçgen varsa linear interpolasyon \\
$\Pi$ varsa kare darbelerle temsil edilir

\noindent (yeniden oluşturma denklemi) Reconstruction eq \\
(İdeal interpolasyon //) ideal interpolation

\sayfa{3}
\elcizimi[0.7]{sekiller/Ders-4-s3-1.png}[3]

\section{Örnek: Lathi; Prob. 6.1-4}
\noindent A signal $x(t)=\mathrm{sinc}^2(5t)$ is sampled (using uniformly spaced impulses) at a rate of:\\
i) 5 Hz, \quad ii) 10 Hz, \quad iii) 20 Hz

\noindent For each of three cases\\
a) Sketch the sampled signal\\
b) Sketch the spectrum of the sampled signal\\
c) Explain whether you can recover the signal $x(t)$ from the sampled signal\\
d) If the sampled signal is passed through an ideal LPF of BW 5 Hz, Sketch the spectrum of the output signal.

\elcizimi[0.85]{sekiller/Ders-4-s3-2.png}[3]

\formulkutu{AT\,\mathrm{sinc}^2(tT)\overset{F}{\longleftrightarrow}A\,\Lambda\!\left(\frac{f}{T}\right)}

\[ x(t)=\mathrm{sinc}^2(5t)\quad\overset{F}{\longleftrightarrow} \]
\noindent $A=\dfrac{1}{T}=\dfrac{1}{5}$

\elcizimi[0.4]{sekiller/Ders-4-s3-3.png}[3]

\noindent \textcircled{\small i} \quad $f_s=5\ \mathrm{Hz}\,,\ T_s=\dfrac{1}{f_s}=\dfrac{1}{5}$ sn
\[ x_\delta(t)=\sum_{n=-\infty}^{\infty}\delta\bigl(t-n\underbrace{T_s}_{1/5}\bigr) \]
\[ x_s(t)=x(t)\,x_\delta(t)=\sum_{n=-\infty}^{\infty}x(nT_s)\,\delta(t-nT_s) \]
\[ X_s(f)=f_s\sum_{n=-\infty}^{\infty}X(f-nf_s) \]

\sayfa{4}
\noindent a)

\elcizimi[0.8]{sekiller/Ders-4-s4-1.png}[4]

\[ \mathrm{sinc}\,k=\frac{\sin(\pi k)}{\pi k}=0 \qquad k\in \mathrm{Z}-\{0\} \]

\elcizimi[0.85]{sekiller/Ders-4-s4-2.png}[4]

\noindent c) No, $f_s=5\ \mathrm{Hz}<2f_m=10\ \mathrm{Hz}$

\noindent d)

\elcizimi[0.65]{sekiller/Ders-4-s4-3.png}[4]

\noindent \textcircled{\small ii} \quad $f_s=10\ \mathrm{Hz}\,,\ T_s=\dfrac{1}{f_s}=\dfrac{1}{10}$ sn
\[ x_\delta(t)=\sum_{n=-\infty}^{\infty}\delta\bigl(t-n\underbrace{T_s}_{1/10}\bigr) \]
\[ x_s(t)=x(t)\,x_\delta(t)=\sum_{n=-\infty}^{\infty}x(nT_s)\,\delta(t-nT_s) \]
\[ X_s(f)=\underset{10}{f_s}\sum_{n=-\infty}^{\infty}X\bigl(f-n\underset{10}{f_s}\bigr) \]

\elcizimi[0.75]{sekiller/Ders-4-s4-4.png}[4]

\sayfa{5}
\elcizimi[0.65]{sekiller/Ders-4-s5-1.png}[5]

\noindent c) Yes, $f_s=10\ \mathrm{Hz}=2f_m=10\ \mathrm{Hz}$

\noindent d)

\elcizimi[0.3]{sekiller/Ders-4-s5-3.png}[5]

\elcizimi[0.35]{sekiller/Ders-4-s5-2.png}[5]

\noindent \textcircled{\small iii} \quad $f_s=20\ \mathrm{Hz}\,,\ T_s=\dfrac{1}{f_s}=\dfrac{1}{20}$ sn
\[ x_\delta(t)=\sum_{n=-\infty}^{\infty}\delta\bigl(t-n\underbrace{T_s}_{1/20}\bigr) \]
\[ x_s(t)=x(t)\,x_\delta(t)=\sum_{n=-\infty}^{\infty}x(nT_s)\,\delta(t-nT_s) \]
\[ X_s(f)=\underset{20}{f_s}\sum_{n=-\infty}^{\infty}X\bigl(f-n\underset{20}{f_s}\bigr) \]

\elcizimi[0.9]{sekiller/Ders-4-s5-4.png}[5]

\elcizimi[0.8]{sekiller/Ders-4-s5-5.png}[5]

\sayfa{6}
\noindent c) Yes, $f_s=20\ \mathrm{Hz}>2f_m=10\ \mathrm{Hz}$

\noindent d)

\elcizimi[0.7]{sekiller/Ders-4-s6-1.png}[6]

\section{Pratik örnekleme}
\noindent (Natural sampling)

\elcizimi[0.9]{sekiller/Ders-4-s6-2.png}[6]

\[ x_p(t)=\sum_{n=-\infty}^{\infty}c_n\,e^{j2\pi nf_0t} \]

\elcizimi[0.5]{sekiller/Ders-4-s6-3.png}[6]

\elcizimi[0.5]{sekiller/Ders-4-s6-4.png}[6]

\[ x_p(t)=\sum_{n=-\infty}^{\infty}c_n\,e^{j2\pi nf_st} \]

\elcizimi[0.3]{sekiller/Ders-4-s6-5.png}[6]

\begin{align*}
c_n &= \frac{1}{T_s}\int_{-T_s/2}^{T_s/2}x_p(t)\,e^{-j2\pi nf_st}\,dt=\frac{1}{T_s}\int_{-T/2}^{T/2}\frac{1}{T}\,e^{-j2\pi nf_st}\,dt\\
&= \frac{1}{T_s}\cdot\frac{1}{T}\cdot\left.\frac{e^{-j2\pi nf_st}}{-j2\pi nf_s}\right|_{-T/2}^{T/2}\\
&= \frac{1}{T_s}\cdot\frac{1}{T}\cdot\kutula{\frac{e^{-j\pi nf_sT}-e^{j\pi nf_sT}}{-j2\pi nf_s}}\\
&= \frac{1}{T_s}\,\frac{\sin(\pi nf_sT)}{\pi nf_sT}=\frac{1}{T_s}\,\mathrm{sinc}(nf_sT)
\end{align*}
\[ c_n=\frac{1}{T_s}\,\mathrm{sinc}(nf_sT) \]
\[ X_p(f)=\sum_{n=-\infty}^{\infty}c_n\,\delta(f-nf_s)=\frac{1}{T_s}\sum_{n=-\infty}^{\infty}\mathrm{sinc}(nf_sT)\,\delta(f-nf_s) \]

\noindent Darbe Genlik mod.\ (DGM) deniliyor ama bu değil!!! sonra \\
Pulse amplitude mod (PAM)

\elcizimi[0.6]{sekiller/Ders-4-s6-6.png}[6]

\[ x_s(t)=x(t)\,x_p(t)=x(t)\sum_{n=-\infty}^{\infty}c_n\,e^{j2\pi nf_st} \]
\[ x_s(t)=\sum_{n=-\infty}^{\infty}c_n\,x(t)\,e^{j2\pi nf_st} \]
\[ \downarrow F \qquad c_n=f_s\cdot\mathrm{sinc}(nf_sT) \]
\[ X_s(f)=\sum_{n=-\infty}^{\infty}c_n\,X(f-nf_s) \]

\sayfa{7}
\elcizimi[0.9]{sekiller/Ders-4-s7-1.png}[7]

\begin{align*}
x_s(t) &\longleftrightarrow X_s(f)=X(f)\otimes X_p(f)=X(f)\otimes\sum_{n=-\infty}^{\infty}c_n\,\delta(f-nf_s)\\
&=\sum_{n=-\infty}^{+\infty}c_n\,X(f-nf_s)
\end{align*}

\section{Örnek: Lathi, Pr. 6.1-3}
\noindent A signal $x(t)$ band limited to $B$ Hz is sampled by a periodic pulse train $p_{T_s}(t)$ made up of a rectangular pulse of width $\frac{1}{8B}$ seconds (centered at the origin) repeating at the Nyquist rate ($2B$ pulses per second). Show that the sampled signal $x_s(t)$ is given by
\[ x_s(t)=\frac{1}{4}\,x(t)+\sum_{n=1}^{\infty}\frac{2}{n\pi}\sin\!\left(\frac{n\pi}{4}\right)x(t)\cos(n\omega_st)\,,\qquad \omega_s=2\pi f_s=4\pi B \]
\noindent Show that the signal $x(t)$ can be recovered by passing $x_s(t)$ through an ideal LPF of BW $B$ Hz and a gain of 4.

\noindent \okunamadi{}

\noindent $f_s=2B$ \\ $T_s=\dfrac{1}{2B}$

\elcizimi[0.65]{sekiller/Ders-4-s7-2.png}[7]

\[ x_s(t)=x(t)\cdot p_{T_s}(t) \]
\[ p_{T_s}(t)=\frac{1}{4}+\sum_{n=1}^{\infty}\frac{2}{n\pi}\sin\!\left(\frac{n\pi}{4}\right)\cos(n\omega_st)\qquad \omega_s=2\pi f_s=4\pi B \]
\noindent böyle olduğunu ispatlamamız gerek
\[ p_{T_s}(t)=a_0+\sum_{n=1}^{\infty}a_n\cos(n\omega_st)+b_n\sin(n\omega_st)\,] \]
\begin{align*}
a_0 &= \frac{1}{T_s}\int_{T_s}x(t)\,dt=\frac{1}{T_0}\int_{-T_s/2}^{T_s/2}x(t)\,dt\\
&= 2B\int_{-1/16B}^{1/16B}1\,dt=2B\cdot\frac{2}{16B}\\
&= \frac{1}{4}
\end{align*}

\sayfa{8}
\noindent $b_n=0$ \quad çünkü fonksiyon çift ($p_{T_s}$ çift $\rightarrow b_n=0$.)
\begin{align*}
a_n &= \frac{2}{T_s}\int_{-T_s/2}^{T_s/2}p_{T_s}(t)\cos(n\omega_st)\,dt=4B\int_{-1/16B}^{1/16B}1\cdot\cos(n\omega_st)\,dt\\
&= 4B\left.\frac{\sin(n\omega_st)}{n\,\underset{4\pi B}{\omega_s}}\right|_{-1/16B}^{1/16B}\\
&= \frac{1}{n\pi}\left[\sin\!\left(n\cdot4\pi B\cdot\frac{1}{16B}\right)-\sin\!\left(n\cdot4\pi B\cdot\frac{-1}{16B}\right)\right]
\end{align*}
\formulkutu{a_n=\frac{2}{n\pi}\sin\!\left(\frac{n\pi}{4}\right)}
\[ a_0=\frac{\pi}{4}\footnotemark\ \rightarrow\ p_{T_s}(t)=\underbrace{\frac{1}{4}}_{a_0}+\underbrace{\sum_{n=1}^{\infty}\frac{2}{n\pi}\sin\!\left(\frac{n\pi}{4}\right)}_{a_n}\cos(n\omega_st) \]
\footnotetext{Claude notu: muhtemelen $a_0=\frac{1}{4}$ olmalı.}
\[ x_s(t)=x(t)\,p_{T_s}(t)=\frac{1}{4}\,x(t)+\sum_{n=1}^{\infty}\frac{2}{n\pi}\sin\!\left(\frac{n\pi}{4}\right)\underbrace{\cos(n\omega_st)}_{e^{j2\pi nf_st}}\,x(t) \]

\noindent b)

\elcizimi[0.45]{sekiller/Ders-4-s8-1.png}[8]

\noindent olsun
\[ X_s(f)=\frac{1}{4}\,X(f)+\frac{1}{2}\,\frac{2}{n\pi}\sin\!\left(\frac{n\pi}{4}\right)\bigl[X(f-nf_s)+X(f+nf_s)\bigr] \]

\elcizimi[0.85]{sekiller/Ders-4-s8-2.png}[8]

\elcizimi[0.5]{sekiller/Ders-4-s8-3.png}[8]

\sayfa{9}
\section{Düz Tepeli örnekleme}
\noindent (Flat top sampling) (FTS)

\elcizimi[0.55]{sekiller/Ders-4-s9-1.png}[9]

\noindent tepeyi takip etmiyor

\elcizimi[0.5]{sekiller/Ders-4-s9-2.png}[9]

\noindent Sample and hold circuit \\ (Örnekle ve tut devresi)

\noindent $x_s(t)$: Darbe Genlik mod (DGM) \\ Pulse amplitude mod

\end{document}
